\section{Model automatu stanów}
\label{sec:automat}

\subsection{Stany i ich kodowanie}
Układ opisano jako automat skończony. Stan kodują trzy zmienne
pamięci: $M_1$ (telefon), $M_2$ (zapamiętany stan ekranu) i $M_3$
(menu). Stan menu otwartego rozdzielono na dwa warianty, ponieważ po
anulowaniu menu ekran musi wrócić do stanu sprzed jego otwarcia
(wymaganie R4). Wymaga to zapamiętania, czy ekran był wcześniej
włączony ($S_{3b}$), czy wygaszony ($S_{3a}$). Pięć stanów
automatu przedstawiono w tabeli~\ref{tab:stany}.

\begin{table}[H]
\centering
\caption{Stany automatu i ich kodowanie.}
\label{tab:stany}
\begin{tabular}{llccc|cc}
\toprule
Stan & Opis & $M_1$ & $M_2$ & $M_3$ & $Q_2$ & $Q_1$ \\
\midrule
$S_0$    & telefon wyłączony                    & 0 & 0 & 0 & 0 & 0 \\
$S_1$    & praca, ekran wygaszony               & 1 & 0 & 0 & 1 & 0 \\
$S_2$    & praca, ekran włączony                & 1 & 1 & 0 & 1 & 1 \\
$S_{3a}$ & menu otwarte (wejście ze stanu $S_1$)& 1 & 0 & 1 & 1 & 1 \\
$S_{3b}$ & menu otwarte (wejście ze stanu $S_2$)& 1 & 1 & 1 & 1 & 1 \\
\bottomrule
\end{tabular}
\end{table}

Pozostałe trzy kombinacje $(M_1M_2M_3)$, to jest $010$, $001$ i $011$,
odpowiadają telefonowi wyłączonemu z włączonym ekranem lub menu i są
niedozwolone (wymaganie R7). W punkcie~\ref{sec:niedozwolone}
pokazano, że program ich nie generuje.

\subsection{Tabela przejść}
Tabela~\ref{tab:przejscia} przedstawia stan następny w zależności
od stanu bieżącego i zdarzenia (K, Ś, D).

\begin{table}[H]
\centering
\caption{Tabela przejść automatu.}
\label{tab:przejscia}
\begin{tabular}{l|ccc}
\toprule
Stan bieżący & K & Ś & D \\
\midrule
$S_0$    & $S_0$    & $S_0$    & $S_2$ \\
$S_1$    & $S_2$    & $S_{3a}$ & $S_0$ \\
$S_2$    & $S_1$    & $S_{3b}$ & $S_0$ \\
$S_{3a}$ & $S_1$    & $S_{3a}$ & $S_0$ \\
$S_{3b}$ & $S_2$    & $S_{3b}$ & $S_0$ \\
\bottomrule
\end{tabular}
\end{table}

Graf przejść pokazano na rysunku~\ref{fig:automat}. Kliknięcie D
z dowolnego stanu z włączonym telefonem prowadzi do $S_0$, dlatego na
rysunku przedstawiono je jedną strzałką wychodzącą z ramki grupującej
stany $S_1$, $S_2$, $S_{3a}$ i $S_{3b}$.

\begin{figure}[H]
\centering
% Diagram automatu stanow -- TikZ
\begin{tikzpicture}[
    >=Stealth,
    node distance=3.6cm,
    st/.style={circle,draw,thick,minimum size=1.15cm,inner sep=1pt,font=\sffamily\small},
    ev/.style={font=\small,fill=white,inner sep=1.5pt}
]
    \node[st] (S0)  at (0,0)    {$S_0$};
    \node[st] (S1)  at (5,1.9)  {$S_1$};
    \node[st] (S2)  at (5,-1.9) {$S_2$};
    \node[st] (S3a) at (10.5,1.9)  {$S_{3a}$};
    \node[st] (S3b) at (10.5,-1.9) {$S_{3b}$};

    % telefon wylaczony
    \path[->,thick]
        (S0) edge[loop above] node[ev] {K, Ś} (S0)
        (S0) edge[bend right=18] node[ev] {D} (S2);

    % przelaczanie ekranu
    \path[->,thick]
        (S1) edge[bend right=18] node[ev] {K} (S2)
        (S2) edge[bend right=18] node[ev] {K} (S1);

    % otwarcie i anulowanie menu
    \path[->,thick]
        (S1)  edge[bend left=18]  node[ev] {Ś} (S3a)
        (S3a) edge[bend left=18]  node[ev] {K} (S1)
        (S2)  edge[bend right=18] node[ev] {Ś} (S3b)
        (S3b) edge[bend right=18] node[ev] {K} (S2)
        (S3a) edge[loop right] node[ev] {Ś} (S3a)
        (S3b) edge[loop right] node[ev] {Ś} (S3b);

    % D z kazdego stanu z wlaczonym telefonem
    \node[draw,dashed,rounded corners,inner sep=0.45cm,
          fit=(S1)(S2)(S3a)(S3b)] (grp) {};
    \node[above=1pt of grp.north west,anchor=south west,font=\small\itshape]
        {stany z włączonym telefonem};
    \draw[->,thick] (grp.west |- S0) -- node[ev,pos=0.5] {D} (S0.east);
\end{tikzpicture}

\caption{Graf przejść automatu obsługi przycisku zasilania.}
\label{fig:automat}
\end{figure}

\subsection{Równania wyjść}
Wyjścia układu wynikają bezpośrednio z kodowania stanów:
\begin{align}
    Q_2 &= M_1, \label{eq:q2}\\
    Q_1 &= M_1\,(M_2 + M_3). \label{eq:q1}
\end{align}
Wyjście $Q_2$ odwzorowuje stan telefonu. Ekran świeci, gdy telefon
działa i ekran jest zapamiętany jako włączony ($M_2$) albo menu jest
otwarte ($M_3$) -- menu wymusza ekran niezależnie od $M_2$.
Czynnik $M_1$ w równaniu~\eqref{eq:q1} gwarantuje spełnienie R7 także
wtedy, gdyby pamięć znalazła się w stanie niedozwolonym.
